\documentclass[11pt]{article}
\usepackage[margin=1in]{geometry}
\usepackage[T1]{fontenc}
\usepackage{lmodern,amsmath,amssymb,booktabs,graphicx}
\usepackage[colorlinks=true,linkcolor=blue,urlcolor=blue]{hyperref}
\setlength{\parindent}{0pt}\setlength{\parskip}{5pt}
\setlength{\emergencystretch}{2em}
\title{RoCE-K: local weak-bias precision and design-branch calibration}
\author{}\date{4 October 2026}
\begin{document}\maketitle

\textbf{Summary.} The linear-model main line now has an all-valid local
length lower bound, derived from exact Bernoulli likelihoods. The construction
also embeds in an i.i.d. randomized population-TATE experiment. A separate
implementation improvement shares the population-composition error event
across design branches, so pre-outcome abstention returns the full-budget
protected target interval. All source/target model assumptions remain explicit.

\section{Accounting and what was held fixed}

There are 4,800 fresh outcome setting-repetitions in 24 local weak-bias cells,
each with 200 repetitions. A further 2,800 repetitions from version 8 are
reanalyzed in pairs under alternative budget/design rules; they are not counted
as new independent datasets. Within the local panel, patterns share random
number streams at each $K$. Exact finite-likelihood calculations are separate
from Monte Carlo repetitions.

The paired panel retains the previous DGP, seeds, feature map, valid-count
assumption, fitted target certificate and point-projection rule. Its legacy
intervals and point estimates reproduce the earlier saved arrays exactly.
Prespecified design policies compare no weight threshold, inflation thresholds
2, 4 and 8, a zero-observed-dispersion reference gate, and threshold 4 plus
that gate. Thresholds are not selected using observed coverage or interval
length. The primary paired panel assumes linear means at every site; the
inference interface also accepts the broader protected model modes.

\section{Budget allocation with a valid design-only fallback}

The old source-assisted ledger spent $.005$ on the target coefficient event,
$.01$ on population composition and $.035$ on conditional outcome/dispersion
events. Its design-fallback branch nevertheless used only $.01$ for target
outcome noise. Increasing both target noise and composition allowances only
after observing which design branch occurred would need additional justification:
the branch may depend on target covariates.

The shared-composition rule instead fixes $.005$ for the coefficient event
and $.0225$ for one composition event across both branches. Conditional on
the full design, the remaining $.0225$ is split across source noise, target
noise and dispersion in the old $.01:.01:.015$ proportions when borrowing;
on abstention it all goes to target noise. Its design-fallback interval is
exactly the previously defined complete-budget protected target-only interval.
This statement is checked separately in every applicable saved repetition.

\begin{center}\small
\resizebox{\linewidth}{!}{\input{experiments/20261004d/budget_table}}
\end{center}

The table reuses the independent-seed confirmation cohort from version 8.
For $p=4,K=2048,h=0$ with weak bias, mean population length changes from
$.1123$ to $.1043$; ordinary target Wald length is $.1227$. Both protected
versions have observed coverage $1.000$, and the Wald reference has $.930$.
The latter's 200-repeat binomial interval still includes $.95$. The new
length is approximately 15\% below the Wald reference in this setting.

The allocation makes conditional source protection slightly wider while
reducing the target composition allowance by more. This is a reallocation
under the same total error guarantee, not a change in the DGP or a use of
true nuisance quantities. At $p=10$ or with heterogeneous target CATE, the
new method remains longer than ordinary Wald inference. The paired improvement
does not constitute another independent-seed replication of the new allocation.

\section{Design variance costs and pre-outcome abstention}

Weight inflation is $n_{ja}\|\gamma_{ja}\|_2^2$ in each arm. The reference
gate evaluates the dispersion certificate at zero observed dispersion using
only retained design factors. It abstains when this reference source radius
is no smaller than the complete-budget target noise radius. This is a
prespecified precision heuristic, not a lower bound on the realized source
interval or a proof that screening is optimal.

\begin{center}\small
\resizebox{\linewidth}{!}{\input{experiments/20261004d/design_budget_table}}
\end{center}

The displayed settings have weak source bias. In the difficult
$n_s=120,p=45,K=64$ setting, the old method retained eligibility but had
length $.3875$. The shared allocation alone gives $.3801$. The reference
gate abstains in every repetition and returns the complete target interval,
length $.3546$, with population/conditional coverage $1.000/.995$.
Threshold 4 also abstains in every repeat. Threshold 8 abstains in only
47.5\% and has average length $.3680$, showing that threshold choice has a
real precision cost. All these prespecified sensitivity results are retained.

At $n_s=200,p=45,K=64$, threshold 4 alone keeps borrowing and gives $.3794$;
the reference gate instead returns the target interval of length $.3539$.
In stable $n_s=1000,p=4,K=256$ designs, neither gate nor these thresholds
abstains. With a prespecified eighth of sources having singular designs,
the shared-budget removal interval has length $.1535$, versus $.1617$ in
version 8 and $.2825$ for protected target-only. Its observed conditional
and population coverage are both $1.000$.

The main text additionally shows that a fixed inflation bound and leverage
bound imply $S_j=O(n_*^{-1})$, $G_j,F_j=O(n_*^{-2})$ and
$d_j^{\max}=O(n_*^{-1})$, where $n_*$ is the smallest retained arm size.
These weaker, observable design bounds already suffice for the conditional
fourth-root rate when adjusted valid proportions stay positive and source
dispersion is local. They do not guarantee attractive finite-sample constants
or remove the target-composition cost.

\section{Local weak-bias experiment}

This panel deliberately isolates the source-inference problem. There are
500 patients in each arm at every site, an exact intercept model, and constant
control mean $.4$. No population-composition enlargement is needed within
this declared model. Source counts are $16,64,256,1024,4096,16384$.

For the mean-matched alternative, target and valid-source treated means are
$.5+\delta_K$, while invalid-source means are $.5-4\delta_K$, where
\[
 \delta_K=\frac{c}{\sqrt{500}\,K^{1/4}},\qquad c\in\{.1,.25\}.
\]
Valid flags are drawn with probability $.8$ and conditioned to contain at
least $\lceil.75K\rceil$ valid sources. These flags are used only to generate
data. Inference knows only the valid-count guarantee. Relative to target,
invalid-source bias is $-5\delta_K$, which vanishes as $K$ grows. All-valid
controls and same-direction shifts of size $5\delta_K$ are also included.

The conditional protected procedure assigns $.0125$ each to source and target
noise and $.025$ to dispersion. The standalone protected target uses $.05$
for outcome noise. Ordinary target Wald and naive source-pooling Wald intervals
use estimated Bernoulli variances and normal critical values. All intervals
target the setting's own effect, including $.1+\delta_K$ in the mean-matched
alternative; the true effect is not incorrectly held at $.1$ across worlds.

\begin{center}\small
\resizebox{\linewidth}{!}{\input{experiments/20261004d/local_table}}
\end{center}

The table uses $c=.25$. At $K=16384$, the invalid-source mean bias is about
$.00494$, versus a single treated-source standard error near $.02236$.
Nonetheless naive pooling has coverage $.020$ because its standard error is
much smaller. The protected interval has coverage $1.000$ and length $.0120$.
The same-direction setting gives naive coverage $.010$ and protected coverage
$1.000$. In the all-valid control naive coverage is $.960$ and protected
coverage $1.000$: paying for weak-bias honesty has a visible length cost even
at favorable data. Small-$K$ protected intervals can remain wider than target
Wald intervals, as the first row shows.

\begin{figure}[ht]\centering
\includegraphics[width=\linewidth]{experiments/20261004d/local_weak_bias.pdf}
\caption{Exact intercept Bernoulli experiment, mean-matched weak deviations
with $c=.25$, 200 repetitions per source count. Both target and every source
contain 1,000 patients. No Gaussian approximation is used by the protected
procedure.}\end{figure}

\section{The corresponding lower bound}

Version 9 defines a local class with maximum source bias bounded by a constant
times $r_{n,K}=\min\{n_s^{-1/2}K^{-1/4},n_t^{-1/2}\}$. A mean-matched
two-point source mixture cancels the first-order likelihood term. Its exact
chi-square divergence starts at order $n_s^2\delta^4$, so $K$ sources cannot
uniformly distinguish shifts on the fourth-root scale. The construction
includes the trusted target information and conditions the source-label prior
to meet the deterministic valid-count guarantee.

The resulting bound is on expected length \emph{at the all-valid null}, for
intervals honest over the local weak-bias class. It is stronger in this respect
than the preceding unrestricted-invalid worst-case bound. An i.i.d. randomized
patient submodel gives the same calculation for population TATE, with covariates
ancillary and constant CATE. For $n_s=n_t=1000$, $K=4096$, $c=.1$ and
valid-flag probability $.8$, the exact audit gives total-variation upper bound
$.0296583$, valid-count tail $3.13\times10^{-15}$, and expected-length lower
bound $.00034403$ at the all-valid null.

The bound's constants were not optimized. It establishes a necessary scale,
not finite-sample optimality of the current interval. With comparable sample
sizes, fixed dimension, stable designs, constant true CATE and $K=o(n^2)$,
the main upper bound has the same order because its extra $n^{-1}$ composition
term is smaller. Heterogeneous CATE, growing target dimension and honest
adaptation over arbitrary nonlinear invalid means remain separate questions.

\section{Validation and remaining work}

The 70 research tests cover exact patient-likelihood identities in both
experiments, count-conditioning and target information, all-valid rate scaling,
legacy inference parity, original valid-count bookkeeping, design-only gate
inputs and exact full-target fallback. Every reported coverage and length is
recomputed from saved endpoints: 140 budget-method rows, 294 paired comparisons
and 96 local-method rows. Earlier code and all results are archived; production
RoCE/ENAR is unchanged.

The next precision question concerns constants and adaptation over the broader
invalid-source class. The new lower bound does not imply that the current
$n_s^{-1/2}$ bounded-invalid certificate floor is optimal at all-valid data.
Additional improvements must preserve the joint error accounting and avoid
using outcome-dependent screening without an accompanying selection argument.
\end{document}
